% CME/Pandoc lettrine support.
% Loaded automatically by scripts/pandoc-cme-xml for LaTeX/PDF output when
% --lettrine is not "none".

\newif\ifcmehaslettrine
\IfFileExists{lettrine.sty}{%
  \usepackage{lettrine}%
  \cmehaslettrinetrue%
}{%
  \cmehaslettrinefalse%
  \PackageWarningNoLine{cme-lettrine}{lettrine.sty not found; CME drop caps will be set inline}%
  \providecommand{\lettrine}[3][]{#2#3}%
}
\IfFileExists{needspace.sty}{\usepackage{needspace}}{\providecommand{\Needspace}[1]{}}

% The lettrine package defaults the rest of the first word to small caps.  CME
% editions keep that text in the body font so decorative modes affect only the
% initial glyph.
\providecommand{\LettrineTextFont}[1]{#1}
\renewcommand{\LettrineTextFont}[1]{#1}

\newcommand{\cmeLettrinePlain}[2]{%
  \lettrine[lines=3,lhang=0.05,nindent=0em]{#1}{#2}%
}

\newcommand{\cmeLettrinePlainAnte}[3]{%
  \ifcmehaslettrine
    \lettrine[lines=3,lhang=0.05,nindent=0em,ante={#1}]{#2}{#3}%
  \else
    #1#2#3%
  \fi
}

\newcommand{\cmeLettrineWithFontHook}[3]{%
  \lettrine[lines=3,lhang=0.05,nindent=0em]{\begingroup#1#2\endgroup}{#3}%
}

\newcommand{\cmeLettrineWithFontHookAnte}[4]{%
  \ifcmehaslettrine
    \lettrine[lines=3,lhang=0.05,nindent=0em,ante={#2}]{\begingroup#1#3\endgroup}{#4}%
  \else
    #2{\begingroup#1#3\endgroup}#4%
  \fi
}

\newcommand{\cmeLettrine}[2]{\cmeLettrinePlain{#1}{#2}}
\newcommand{\cmeLettrineAnte}[3]{\cmeLettrinePlainAnte{#1}{#2}{#3}}
